\subsection{The ring W(A) of Witt vectors}

Let A be a ring. If \(\boldsymbol{a} = (a_n)_{n \in \mathbb{N}}\) and \(\boldsymbol{b} = (b_n)_{n \in \mathbb{N}}\) are elements of \(\mathbf{A}^{\mathbf{N}}\), we will denote by \(S_{\mathbf{A}}(\boldsymbol{a}, \boldsymbol{b})\) (resp. \(P_{\mathbf{A}}(\boldsymbol{a}, \boldsymbol{b})\), respectively \(I_{\mathbf{A}}(\boldsymbol{a})\)), or simply by \(S(\boldsymbol{a}, \boldsymbol{b})\) (resp. \(P(\boldsymbol{a}, \boldsymbol{b})\), respectively \(I(\boldsymbol{a})\)), the sequence \((S_n(a_0, ..., a_n; b_0, \ldots, b_n))_{n \in \mathbb{N}}\) (resp. \((P_n(a_0, ..., a_n; b_0, \ldots, b_n))_{n \in \mathbb{N}}\), respectively \((I_n(a_0, ..., a_n))_{n \in \mathbb{N}}\)). By substituting \(a_n\) for \(X_n\) and \(b_n\) for \(Y_n\), for all \(n \in \mathbb{N}\), in formulas (12), (13), and (14), we obtain the equalities

\begin{enumerate}
\def\labelenumi{(\arabic{enumi})}
\setcounter{enumi}{15}
\tightlist
\item
\end{enumerate}

\[
\Phi_ {\mathrm{A}} (\mathrm{S} _ {\mathrm{A}} (a, b)) = \Phi_ {\mathrm{A}} (a) + \Phi_ {\mathrm{A}} (b)\tag{17}
\]

\[
\Phi_ {\mathrm{A}} (\mathrm{P} _ {\mathrm{A}} (a, b)) = \Phi_ {\mathrm{A}} (a) \Phi_ {\mathrm{A}} (b)\tag{18}
\]

\[
\Phi_ {\mathbf {A}} \left(\mathrm{I} _ {\mathbf {A}} (\boldsymbol {a})\right) = - \Phi_ {\mathbf {A}} (\boldsymbol {a}).
\]

We shall denote W(A) the set A \(^{N}\) endowed with the composition laws S \(_{A}\) and P \(_{A}\).

Let \(\rho : \mathbf{B} \to \mathbf{A}\) be a ring homomorphism. We shall denote by \(\rho^{\mathbf{N}}\), or again by \(\mathbf{W}(\rho)\), the map from \(\mathbf{B}^{\mathbf{N}}\) into \(\mathbf{A}^{\mathbf{N}}\) that sends the element \(b = (b_n)_{n \in \mathbb{N}}\) of \(\mathbf{B}^{\mathbf{N}}\) to \((\rho(b_n))_{n \in \mathbb{N}}\). It follows immediately from the definitions that we have

\begin{enumerate}
\def\labelenumi{(\arabic{enumi})}
\setcounter{enumi}{18}
\tightlist
\item
\end{enumerate}

\[
\mathrm{W} (\rho) \circ \mathrm{S} _ {\mathrm{B}} = \mathrm{S} _ {\mathrm{A}} \circ (\mathrm{W} (\rho) \times \mathrm{W} (\rho))\tag{20}
\]

\[
\mathrm{W} (\rho) \circ \mathrm{P} _ {\mathrm{B}} = \mathrm{P} _ {\mathrm{A}} \circ (\mathrm{W} (\rho) \times \mathrm{W} (\rho))\tag{21}
\]

\[
\mathrm{W} (\rho) \circ \mathrm{I} _ {\mathrm{B}} = \mathrm{I} _ {\mathrm{A}} \circ \mathrm{W} (\rho)\tag{22}
\]

\[
\rho^ {\mathbf {N}} \circ \Phi_ {\mathrm{B}} = \Phi_ {\mathrm{A}} \circ \mathrm{W} (\rho).
\]

Lemma 3. --- Let A be a ring. There exists a surjective ring homomorphism \(\rho: \mathbf{B} \to \mathbf{A}\), where \(\mathbf{B}\) is a ring satisfying the following conditions: p is not a zero divisor in \(\mathbf{B}\), and there exists an endomorphism \(\sigma\) of \(\mathbf{B}\) such that \(\sigma(b) \equiv b^p\) mod. \(p\) . \(\mathbf{B}\) for every \(b \in \mathbf{B}\).

Indeed, it suffices to set \(\mathbf{B} = \mathbf{Z}[(X_{a})_{a \in A}]\), to take for \(\sigma\) the endomorphism of \(\mathbf{B}\) defined by \(\sigma(X_{a}) = X_{a}^{p}\) for every \(a \in A\), and to take for \(\rho\) the homomorphism of \(\mathbf{B}\) into A defined by \(\rho(X_{a}) = a\) for every \(a \in A\).

THEOREM 1. --- a) Let A be a (commutative) ring. Equipped with the addition \(S_A\) and the multiplication \(P_A\), \(W(A)\) is a (commutative) ring. The identity element for addition is the sequence \(0_A\), all of whose terms are zero; the identity element for multiplication is the sequence \(1_A\), all of whose terms are zero except the one of index 0, which equals \(1_A\). The additive inverse of an element \(a\) of \(W(A)\) is \(I_A(a)\).

\begin{enumerate}
\def\labelenumi{\alph{enumi})}
\setcounter{enumi}{1}
\item
  Let \(\rho: \mathbf{B} \to \mathbf{A}\) be a ring homomorphism. Then \(\mathbf{W}(\rho): \mathbf{W}(\mathbf{B}) \to \mathbf{W}(\mathbf{A})\) is a ring homomorphism.
\item
  Let A be a ring. The map \(\Phi_{\mathbf{A}}\) is a ring homomorphism from \(W(A)\) into the product ring \(\mathbf{A}^{\mathbf{N}}\). In particular, for every \(n \in \mathbf{N}\), the map \(\Phi_{n}: a \mapsto \Phi_{n}(a_{0}, \ldots, a_{n})\) is a ring homomorphism from W(A) to A.
\end{enumerate}

Given formulas (16), (17), (19), and (20), it suffices to prove assertion \(a\). Let \(\rho: \mathbf{B} \to \mathbf{A}\) be a ring homomorphism satisfying the conditions of Lemma 3. Let \(\mathbf{B}'\) be the subring of \(\mathbf{B}^{\mathbf{N}}\) consisting of the elements \((b_n)_{n \in \mathbf{N}}\) such that \(\sigma(b_n) \equiv b_{n+1} \pmod{p^{n+1}\mathbf{B}}\) for all \(n \in \mathbf{N}\). By Proposition 2 of no. 2, \(\Phi_{\mathbf{B}}\) induces a bijection \(\Phi_{\mathbf{B}}'\) of W(B) onto \(\mathbf{B}'\). In view of formulas (16) to (18) and the relations \(\Phi_n(\mathbf{0}_{\mathbf{B}}) = 0\) and \(\Phi_n(\mathbf{1}_{\mathbf{B}}) = \mathbf{1}_{\mathbf{B}}\) (\(n \in \mathbf{N}\)), one sees by transport of structure that W(B) is a ring, with identity element \(\mathbf{0}_{\mathbf{B}}\) for addition, \(\mathbf{1}_{\mathbf{B}}\) for multiplication, and opposite of b equal to \(\mathbf{I}_{\mathbf{B}}(b)\).

The map W(ρ): W(B) → W(A) is surjective. By formulas (19) and (20), the equivalence relation R on W(B) associated with the map W(ρ) is compatible with the ring structure of W(B). Since W(ρ) induces a bijection Ψ from the quotient ring W(B)/R onto W(A), compatible with the laws of addition and multiplication, assertion a) follows by transport of structure.

DEFINITION 1. --- Let A be a ring. The ring \(W(A)\) is called the ring of Witt vectors with coefficients in A.

For \(\boldsymbol{a}\) in \(W(A)\) and \(n\) in \(\mathbf{N}\), the element \(\Phi_n(\boldsymbol{a}) = \Phi_n(a_0, \ldots, a_n)\) is sometimes called the ghost component of index \(n\) of \(\boldsymbol{a}\).

Remark. --- We take up again the notations of the remark of no. 3. Let A be a ring. If \(a\) and \(b\) are elements of \(\mathbf{A}^{\mathbf{J}}\) and \(r = (r_{j})_{j\in\mathbf{J}}\) is an element of \(\mathcal{A}^{\mathbf{J}}\), denote by \(r_{\mathrm{A}}(a,b)\) the element \((r_{j}(a,b))_{j\in\mathbf{J}}\) of \(\mathbf{A}^{\mathbf{J}}\). Denote by \(U(A)\) the set \(\mathbf{A}^{\mathbf{J}}\) endowed with the composition laws \(s_{\mathrm{A}}\) and \(p_{\mathrm{A}}\). One can show (p. 52, exerc. 35) that, endowed with the addition \(s_{\mathrm{A}}\) and the multiplication \(p_{\mathrm{A}}\), \(U(A)\) is a (commutative) ring; it is called the universal Witt ring of A. The identity element for addition is the element of \(U(A)\) all of whose components are zero; the identity element for multiplication is the element of \(U(A)\) all of whose components are zero except the one of index 1, which equals \(1_{\mathrm{A}}\); the opposite of an element \(a\) of \(U(A)\) is \(i_{\mathrm{A}}(a)\). The map \(\varphi_{\mathrm{A}}\) is a ring homomorphism from \(U(A)\) to the product ring \(\mathbf{A}^{\mathbf{J}}\).

Let \(\rho: \mathbf{B} \to \mathbf{A}\) be a ring homomorphism; denote by \(\mathrm{U}(\rho)\) the map from \(\mathbf{B}^{\mathbf{J}}\) into \(\mathbf{A}^{\mathbf{J}}\) that sends the element \((b_j)_{j \in \mathbf{J}}\) of \(\mathbf{B}^{\mathbf{J}}\) to the element \((\rho(b_j))_{j \in \mathbf{J}}\) of \(\mathbf{A}^{\mathbf{J}}\). One can show (loc. cit.) that \(\mathrm{U}(\rho)\) is a ring homomorphism from \(\mathrm{U}(\mathrm{B})\) into \(\mathrm{U}(\mathrm{A})\).
% semantic-fragments: legacy-input-seam
\relax{}
